\workbookchapter{模型安全}

\begin{exercises}
\begin{exercise} 为一个仅回答问题的模型和一个可修改业务记录的模型分别列出资产、不可接受后果和新增信任边界。

\begin{solution}
问答模型的资产含输入文档、隐私、回答可信度和来源；不可接受后果包括泄露或高损失误导。增加写权限后还涉及账户、库存、订单与业务状态，新增模型提案到授权网关、网关到业务系统、批准到执行的信任边界。控制需从内容检查扩展至对象权限、参数约束、幂等与回执核验，不能仅提高拒答率。
\end{solution}
\end{exercise}

\begin{exercise} 推导期望损失公式在互斥场景与可加损失分量两种定义下的含义，构造攻击路径重叠导致重复计数的例子。

\begin{solution}
互斥完备场景$S_i$下，$E[L]=\sum_iP(S_i)E[L\mid S_i]$。若定义可加损失分量$L=\sum_jL_j$，线性期望给$E[L]=\sum_jE[L_j]$，不需独立。若两条攻击路径均导致同一次100单位泄露且各概率0.1、完全重叠，分别算0.1乘100再相加得到20而真实为10；路径不是互斥场景，也不是两个不同损失分量，须去重或按最终事件划分。
\end{solution}
\end{exercise}

\begin{exercise}\optional 使用不同单调编码展示序数风险乘积排序可改变，说明为什么它不能解释为真实风险倍数。

\begin{solution}
设风险A等级$(1,4)$、B$(2,2)$，编码1、2、3、4时乘积均4；采用保持顺序的编码$(1,3,4,5)$时A为5、B为9，B高；编码$(1,2,3,10)$时A为10、B为4，A高。等级只规定顺序，任意单调重标度可改变乘积关系，因此不能解释为概率乘真实损失或风险倍数。
\end{solution}
\end{exercise}

\begin{exercise} 两层控制的漏放概率分别为0.2和0.3，给出联合漏放率的可行范围，并构造独立与相关失效情形。

\begin{solution}
设漏放事件A、B，Fréchet界为$\max(0,0.2+0.3-1)\le P(A\cap B)\le\min(0.2,0.3)$，即0至0.2。独立时为0.06；A包含于B时为0.2；两者互斥时为0。不能默认层数增加就乘法降低风险，应测条件漏放$P(B\mid A)$和共同失败原因。
\end{solution}
\end{exercise}

\begin{exercise} 对内部报告外发场景，分别指出模型行为控制、授权控制、凭据范围与输出渲染能够阻断的路径。

\begin{solution}
模型行为控制可减少建议外发的概率；授权控制检查任务目的、收件对象和数据级别，阻止即便格式合法的越权发送；短期最小凭据限制可触达资源和端点；渲染控制防止把工具文本中的主动资源或链接自动执行/加载。它们保护不同路径，内容分类器不能代替确定性授权，工具成功也不证明数据允许披露。
\end{solution}
\end{exercise}

\begin{exercise} 设计一个规范化资源的例子，说明先检查字符串、后解析资源可能怎样破坏授权对应关系；只讨论防御契约，不执行攻击。

\begin{solution}
防御例可用对象别名或大小写、编码差异：检查时把别名视为公开对象，解析器最终映射到受限对象。应先用唯一规定规则规范化并解析为稳定资源ID，再按该ID和当前租户授权，执行时使用同一解析结果及版本。若重定向或资源映射可变化，重新校验最终对象；拒绝歧义输入，不维护两套不一致的路径解释。
\end{solution}
\end{exercise}

\begin{exercise} 为需要批准的动作设计摘要字段，说明参数、权限或资源版本改变后如何判定原批准是否仍有效。

\begin{solution}
摘要字段包括主体/租户、动作类型、工具版本、最终资源ID、规范参数、目的地址、数据级别、对象/权限版本和有效期。参数或资源版本变化后重算摘要并检查批准绑定；权限撤销即使摘要不变也应阻止执行。批准只在其明确范围及期限内有效，不是给予模型后续任意操作的令牌。
\end{solution}
\end{exercise}

\begin{exercise} 工具已经执行但响应超时，设计执行账本和对账状态机，说明为何使用新幂等键重试可能错误。

\begin{solution}
账本按租户和稳定幂等键唯一认领，绑定动作摘要，状态记录发送中、未知、确认成功或明确失败。超时后按原键查询工具结果，确认成功则补回执；不能确认则保持未知并人工对账。新键使服务端把同一意图当新动作，可能再次扣款或发信。若远端没有幂等或可靠查询，只能停止自动重发，不能由本地账本宣称恰好一次。
\end{solution}
\end{exercise}

\begin{exercise} 比较哈希一致、签名可信、加载成功和行为评测通过分别能够支持什么结论。

\begin{solution}
哈希一致说明与给定字节基准相同，前提是基准可信；签名说明受信密钥对某制品背书，但不保证行为安全；加载成功说明格式与部分环境兼容；行为评测通过只支持指定样本和协议下的观察。四者均不独立证明生产安全，需要制品身份、可信来源、隔离加载和任务回归的联合证据。
\end{solution}
\end{exercise}

\begin{exercise}\optional 设计样本到数据快照、模型和部署的谱系，说明删除原始数据后仍需分析哪些派生资产。

\begin{solution}
谱系可为样本ID到清洗/切片修订、训练快照、优化运行、模型/适配器哈希、索引/缓存、部署修订。删除原文后还需检查副本、备份、向量、摘要、日志、评测集及已训练参数等派生资产；不同资产删除语义不同。追踪用于识别影响范围，不能仅删除一条数据库记录就声称模型已遗忘该信息。
\end{solution}
\end{exercise}

\begin{exercise} 根据不同危险请求基率计算拦截精确率，解释为什么同一过滤器在不同业务中审核负担不同。

\begin{solution}
题目未指定参数，取危险召回$r=0.9$、正常误拦$f=0.02$，基率$\pi$时拦截精确率为 $\frac{r\pi}{[r\pi+f(1-\pi)]}$。$\pi=0.01$得$\frac{0.009}{0.0288}=31.25\%$；$\pi=0.1$得$\frac{0.09}{0.108}=83.33\%$；$\pi=0.001$仅约4.31\%。相同过滤器在低基率下大多数拦截可能是假阳性，审核负担应按请求量与误拦率共同估计。
\end{solution}
\end{exercise}

\begin{exercise} 从零失败概率推导单侧置信上界；若要在独立同分布前提下将95\%上界降至1\%以下，推导最小样本数。

\begin{solution}
若独立同分布失败概率q，n次零失败概率$(1-q)^n$；令其等于$\alpha$反解单侧上界$q_U=1-\alpha^{\frac{1}{n}}$。要求95\%上界严格小于0.01，须$n>\frac{\log0.05}{\log0.99}\approx298.073$，最小整数299。零失败不代表q为0，相关样本或变化分布会破坏该二项界的前提。
\end{solution}
\end{exercise}

\begin{exercise} 为多语言安全评测设计正常、风险和无正常样本三类切片，明确定义漏放、误拦及控制覆盖的分母。

\begin{solution}
每语言分别准备正常、风险及分层场景；漏放分母为该语言真实危险请求数，误拦分母为真实正常请求数，控制覆盖分母为预先列出的应受控路径或场景数并注明加权方式。若某切片只有危险样本，误拦率未定义，不得填0；只有正常样本则危险漏放未定义。总体聚合须按目标分布权重，不能用多语言题量代替真实流量比例。
\end{solution}
\end{exercise}

\begin{exercise} 将构造事件扩展为结果未知而非已成功的工具调用，说明调查、遏制与恢复流程中哪些判断必须改变。

\begin{solution}
把10:01成功改为发送后超时，则10:03告警不能证明已泄露，也不能证明未执行。先停用发送并撤销凭据限制扩大，同时保全动作键、收件信息及工具端日志；按原键查询和人工对账确认影响。检测时延若无已确认执行时刻应报告从首次可疑行为或告警起的明确替代口径，不能继续写执行到告警2分钟。恢复前先处理未知动作，防止恢复后盲重放造成真正泄露。
\end{solution}
\end{exercise}

\end{exercises}
